\begin{definition}[Taylor polynomial of degree $k$]
Let $k\in\mathbb{N}$, let $u\colon\mathrm{Plane}\to\mathbb{R}$ and let
$x_0\in\mathrm{Plane}$.  The \emph{Taylor polynomial of $u$ of degree $k$ at $x_0$} is the
function $\mathrm{TaylorPoly}(k,u,x_0)\colon\mathrm{Plane}\to\mathbb{R}$ defined by
\[
  \mathrm{TaylorPoly}(k,u,x_0)(z)
  :=\sum_{j=0}^{k}\frac{1}{j!}\,D^{j}u(x_0)\bigl[z-x_0,\ldots,z-x_0\bigr],
\]
where $D^{j}u(x_0)$ is the $j$-th iterated Fr\'echet derivative of $u$ at $x_0$, a continuous
$j$-linear form on $\mathrm{Plane}$, and the bracket contains $j$ copies of $z-x_0$.  For
$j=0$ the term is $D^{0}u(x_0)[\;]=u(x_0)$.

This is the classical Taylor expansion of $u$ about $x_0$ truncated at order $k$, and it is
a polynomial function of degree at most $k$ on $\mathrm{Plane}$.  The definition is purely
formal: it makes sense for an arbitrary $u$, and no regularity of $u$ is assumed.  Only the
\emph{diagonal} evaluations $D^{j}u(x_0)[w,\ldots,w]$ of the iterated derivatives enter, in
accordance with the convention of this development that higher-order derivatives are always
carried through their diagonal evaluations.  We write $P$ for
$\mathrm{TaylorPoly}(k,u,x_0)$ when $k$, $u$ and $x_0$ are clear from the context, matching
the paper's notation.
\end{definition}
